\chapter{Introdución}
\label{chap:introducion}

\lettrine{P}{rimer} capítulo da memoria, onde xeralmente se exporán as
liñas mestras do traballo, os obxectivos, etc. Incluimos un par de
exemplos de citas~\cite{ErlangBook,ErlangWebBook} e de referencias
internas (sección \ref{sec:mostra}, páxina \pageref{sec:mostra}).

\lettrine{L}a inteligencia artificial ...?


\section{Contexto y motivación}
\label{sec:contexto_y_motivacion}

Los algoritmos de aprendizaje máquina son técnicas matemáticas que escapan de las fronteras de lo que es entendible a simple vista por ¿?una persona cualquiera.¿?
Desde sus raices más puramente matemáticas hasta su implementación tanto hardware como software. 

Esto ya era cierto en los tiempos de los primeros algoritmos de aprendizaje ¿?por refuerzo, árboles o qué¿?,
pero con el surgir de los modelos de aprendizaje profundo, esta realidad es más palpable que nunca.

Bajo estas premisas es que se siente indispensable ¿? proporcionar a la gente los materiales y métodos para que puedan comprender lo mejor posible esta tecnología cada vez más omnipresente¿?.
¿? Y a de más el aprendizaje es mejor si es interactivo ¿?



\subsection{Definición del Problema}

Hoy en día las principales formas de aprender cómo funcionan los modelos de inteligencia artificial (sobre todo los clásicos, pero también los modernos) es a traves de los siempre confiables libros y universidad, las más ancestrales fuentes de conocimiento.
Entre los formatos más actuales también encontramos los vídeos divulgativos en plataformas online como Youtube,
dónde la posibilidad de que todo el mundo comparta su contenido ha dado lugar a ¿? muchas oportunidades de ver las interpretaciones visuales de otras personas¿?.

Sin embargo ¿? aún no existe ninguna herramienta pedagógica que mermita interactuar en tiempo real con los algorimos.¿?
¿?Programándolos se pueden ver sus resultados, su evolución en una gráfica,
pero no palpar los que está sucediendo en las entrañas de la bestia mientras mastica el equivalente a su comida: los datos,
tanto en entrenamiento como en inferencia¿? ((.......))


\subsection{Objetivos}

((
Crear la aplicación
- paso 1, hacerla
- paso 2, exportarla

Publicarla para que todo el mundo la pueda disfrutar
))

((creo qe hay q se rmás concreto))

Los objetivos que se pantean en este trabajo comienzan por plantearse qué algoritmos de IA son los más interesantes para ¿?explicar¿?.
Esta es, en última instancia, una cuestión subjetiva,
por lo que en este trabajo se abordó esta decisión mirando ¿?de cara¿? a dos de los mayores paradigmas de la IA: la simbólica y el deeplearning.
((no se si decir de forma pragmática y decir que tenía en mente las redes neuronales, y que los árboles se parecían lo suficiente, o poner un poco de fantasía de "los árboles, uno los primeros amigos del hombre en cuánto a lo que hacer que una máquina aprenda a jugar al ajedrez..." (lo 

De ambos se decidió que los mejores candidatos para protagonizar una explicacición que busque asentar los fundamentos eran las versiones más fundamentales de cada uno.
Del lado de la IA simbólica está el árbol de decisión ((planteado por bla bla bla en blabal)), conocido por ser ... white box
Y en el deeplearning, el ladrillo fundacional, la red neuronal ((lo mismo de histoia...))


\subsection{Alcance y Limitaciones}




\subsection{Metodología de Trabajo}



\subsection{Estructura de la Memoria}

